\chapter{Introdução}

\section{Contextualização}
\section{Motivação}
\section{Problema de Pesquisa}
\section{Objetivos}
\subsection{Objetivo Geral}
\subsection{Objetivos Específicos}
\section{Delimitação do Escopo}
\section{Estrutura do Trabalho}

\chapter{Fundamentos de Armazenamento de Dados em DNA e Canais de Erro}

\section{DNA como Meio de Armazenamento}
\section{Codificação Digital para DNA}
\section{Tipos de Erro no Canal de DNA}

\chapter{Códigos LT e DNA Fountain}

\section{Códigos Fountain e Códigos LT}
\section{Distribuição Ideal do Soliton}
\section{Distribuição Robusta do Soliton}
\section{DNA Fountain}
\subsection{Pipeline de Codificação}
\subsection{Estrutura de uma Droplet}
\subsection{Decodificação}

\chapter{Modelo Matemático do Sistema}

\section{Representação em Grafo Bipartido}
\section{Formalização do Peeling Decoder}
\section{Overhead e Condições de Recuperação}
\section{Modelagem do Canal Probabilístico de Apagamento}
\section{Hipótese do Trabalho}

\chapter{Metodologia de Simulação em Software}

\section{Simulador de Referência}
\section{Geração de Instâncias}
\section{Definição dos Parâmetros}
\section{Procedimento Experimental}
\section{Métricas de Avaliação}

\chapter{Resultados e Discussão}

\section{Taxa de Recuperação em Função do Overhead}
\section{Efeito da Taxa de Apagamento}
\section{Efeito dos Parâmetros da Distribuição Robusta do Soliton}
\section{Comportamento do Ripple}
\section{Discussão dos Resultados}

\chapter{Conclusão}

\section{Conclusões}
\section{Limitações do Trabalho}
\section{Trabalhos Futuros}